\begingroup

\chapter*{Introduction: Spectral Contexts and the Infinitesimal Foundations of Variational Geometry}
\phantomsection
\addcontentsline{toc}{chapter}{Introduction}

\setcounter{section}{0}
\renewcommand{\thesection}{\arabic{section}}
\makeatletter
\@ifundefined{theHsection}{}{%
    \renewcommand{\theHsection}{intro.\arabic{section}}%
}
\makeatother


\section{From variational geometry to the spectral problem}
\label{sec:intro-programme}

Variational geometry studies fields through their dependence on parameters,
their infinitesimal variation, and the differential equations that constrain
them.  In classical geometry this leads to jet spaces, formal integrability,
variational sequences, and the variational bicomplex
\cite{Kock1980,Marvan1986,Vitolo2000,GiachettaMangiarottiSardanashvily2000,PuglieseSparanoVitagliano2023}.
In mathematical physics these structures meet symplectic and multisymplectic
geometry, gauge symmetry, cohomological information, and higher geometric
structures
\cite{Rogers2012,PTVV2013,Calaque2013,FiorenzaSatiSchreiber2013,FiorenzaRogersSchreiber2013,AlfonsiYoung2023}.
A homotopical formulation of variational geometry should therefore retain the
deformation and coherence data carried by these interactions while still
supporting the local and infinitesimal constructions on which differential
geometry depends.

Spectral algebraic geometry provides a setting in which such deformation data
are intrinsic, but it also introduces the structural problem that organizes
this dissertation.  Its nonlinear affine geometry is built from connective
commutative ring spectra, whereas its natural coefficient categories are
fully stable module categories
\cite{LurieHigherAlgebra,LurieSAG}.  The connective layer governs spectral
schemes, relative affine spaces, and nonlinear context extension.  The stable
layer carries cotangent complexes, mapping spectra, fibres, cofibres,
suspensions, and totalizations.  A spectral approach to variational geometry
must use both layers without identifying them, and it must make their
interaction coherent under change of geometric base.

This dissertation is the first stage of a longer programme aimed at
developing such a framework.  The eventual goal is a spectral-algebraic
formulation of variational geometry suitable for foundational questions in
mathematical physics, with higher gauge theory and derived field theory as
principal sources of motivation and intended applications.  The present work
develops the
geometric and infinitesimal infrastructure needed before a genuinely
variational theory can be formulated.  It reaches spectral infinitesimal
variation, formal reconstruction, de Rham descent, crystalline transport,
infinite jets, and a formal differential-equation calculus.  Lagrangian
structures, Euler--Lagrange equations, variational complexes, and a spectral
theory of higher gauge fields are not constructed here and belong to later
stages of the programme.

The first problem is organizational: the connective nonlinear
geometry and the stable linear coefficients must be assembled into one
indexed system.  We seek a geometry in which change of base transports
coefficients, local data satisfy descent, relative algebra determines geometry
over a fixed base, and infinitesimal extensions encode variation.  The
language of categorical logic is useful for expressing this organization.  A
spectral scheme is read as a \emph{context}, pullback as \emph{substitution},
and relative spectrum as \emph{comprehension} or context extension
\cite{JacobsCategoricalLogic}.  These terms describe the indexed role of the
constructions; the constructions themselves remain spectral-geometric and
higher-categorical.

The guiding question of the dissertation is:

\begin{quote}
How much of the infinitesimal and jet-theoretic foundation of variational
geometry can be constructed from connective spectral contexts and their fully
stable quasi-coherent coefficients?
\end{quote}

The answer developed here is that these two layers already support a coherent
infinitesimal and infinite-jet calculus.  The construction keeps fibrewise
stable differentiation distinct from nonlinear infinitesimal geometry,
connects them through the cotangent classifier, and then passes from finite
infinitesimal variation to reconstruction, de Rham localization, effective
descent, and crystalline jet transport.  On specified complete fixed-centre
domains, the nonlinear augmentation is reconstructed from shifted-cotangent
data together with the coalgebra coherence generated by structured
square-zero formation; separately, Postnikov continuity and nilpotent
excision define a centre-free reconstruction localization.
Making finite infinitesimal variation invisible produces a de Rham reflector
and an effective quotient.  Dependent adjoints along that quotient then
produce the formal-disk and infinite-jet operators from which crystalline
transport and formal differential equations are obtained.

\section{The connective--stable architecture of infinitesimal variation}
\label{sec:intro-spectral-contexts}

The programme requires nonlinear spectral geometry and stable linear
coefficients to vary together without being collapsed into a single kind of
object.  We therefore organize spectral contexts, quasi-coherent
coefficients, and change of base into one indexed geometry, with relative
spectrum providing the return from connective coefficient algebra to
geometric context extension.

The fixed affine background is
\[
    \Aff_{\Sp}
    =
    \bigl(\CAlg(\Sp)^{\mathrm{cn}}\bigr)^{\op}.
\]
For an affine context \(X=\Spec(A)\), the coefficient fibre is
\[
    \QCoh_{\mathrm{aff}}(X)=\Mod_A,
\]
and a morphism
\[
    f:\Spec(B)\longrightarrow \Spec(A)
\]
represented by \(A\to B\) acts on coefficients by extension of scalars,
\[
    f^*=-\otimes_A B:\Mod_A\longrightarrow\Mod_B.
\]
Spectral Zariski descent globalizes this assignment to
\[
    \QCoh:
    \bigl(\operatorname{Sch}^{\mathrm{sp}}\bigr)^{\op}
    \longrightarrow
    \CAlg(\Pr^L_{\mathrm{st}}),
\]
so every spectral context carries a stable symmetric-monoidal category of
quasi-coherent coefficients, and every change of context carries a coherent
strong symmetric-monoidal substitution functor
\cite{LurieHigherAlgebra,LurieSAG}.

This coefficient doctrine remains connected to nonlinear geometry through
relative spectrum.  For a spectral scheme \(X\), relative spectrum identifies
connective quasi-coherent commutative algebras with affine morphisms over
\(X\):
\[
    \operatorname{QCAlg}(X)^{\mathrm{cn},\op}
    \simeq
    \bigl(\operatorname{Sch}^{\mathrm{sp}}_{/X}\bigr)_{\mathrm{aff}}.
\]
Thus, for
\(\mathcal A\in\operatorname{QCAlg}(X)^{\mathrm{cn}}\), the context extension
\[
    \Spec_X(\mathcal A)\longrightarrow X
\]
represents algebra terms after substitution.  This is the comprehension
principle used throughout the dissertation.  Crucially, comprehension returns \emph{connective} coefficient algebra to
geometry; it does not identify the geometric and stable layers.

The connective--stable distinction remains active after comprehension.
Geometry is built from connective commutative algebra, whereas
cotangent complexes, mapping spectra, stabilizations, fibres, cofibres, and
totalizations live in fully stable coefficient categories.  When
cocommutative coalgebras appear below,
\(\operatorname{CoCom}(\mathcal L_X)\) denotes the infinity-category of
cocommutative coalgebra objects internal to the stable symmetric-monoidal
fibre \(\mathcal L_X\); stability is asserted for the ambient coefficient
fibre, not for the coalgebra category itself.

With this indexed background in place, differentiation presents two different
problems.  The first asks how coefficients vary \emph{inside} a stable fibre.
The second asks how the nonlinear \emph{context itself} varies
infinitesimally.  The dissertation constructs these two mechanisms
independently and compares them only through a controlled cotangent
interface.

For the fibrewise linear problem, one has affinely the tangent equivalence
\[
    \operatorname{Stab}\bigl(\CAlg(\Sp)_{/B}\bigr)
    \simeq
    \Mod_B.
\]
For a general spectral scheme \(X\), however, the descended stable fibre
\(\QCoh(X)\) is used directly; no global identification of this fibre with a
stabilization of a global nonlinear category is required or asserted.  Writing
\(\mathcal L_X:=\QCoh(X)\), the forgetful functor from cocommutative
coalgebras admits a cofree right adjoint,
\[
    U_X:
    \operatorname{CoCom}(\mathcal L_X)
    \rightleftarrows
    \mathcal L_X:C_X,
    \qquad
    U_X\dashv C_X.
\]
The composite
\[
    !_X=U_XC_X
\]
is the fibrewise differential exponential.  Its canonical coalgebra,
primitive, and bialgebra structures produce the codereliction and deriving
transformation.  After passage to the homotopy category, these maps recover
the corresponding differential-storage structure in the sense of
differential categories
\cite{Benton1995LNL,BluteCockettSeely2006Differential,EhrhardRegnier2003Differential,Lemay2021Coderelictions}.
The existence and behaviour of the spectral cofree coalgebra construction is
closely related to recent work on spectral
\(\mathbb E_\infty\)-coalgebras
\cite{BachmannBurklund2026Coalgebras}.

Change of geometric base does not in general preserve this exponential by an
equivalence.  Instead, for \(f:Y\to X\), cofreedom and pullback construct the
canonical oplax comparison
\[
    \beta_{f,M}:f^*(!_XM)\longrightarrow !_Y(f^*M).
\]
The resulting fibrewise calculus is coherently compatible with substitution
in the oplax sense actually supplied by the construction; no general
equivalence is asserted.

The geometric differentiation problem is controlled by the cotangent
classifier.  For a relative context \(X\to S\), affine cotangent complexes
glue by spectral Zariski descent to a canonically determined object
\[
    L_{X/S}\in\QCoh(X).
\]
Global derivation terms are represented by
\[
    \operatorname{Der}_S(X;M)
    :=
    \Map_{\QCoh(X)}(L_{X/S},M),
\]
recovering on compatible affine charts the usual spectral cotangent
representation of derivations.  Cotangent base change and transitivity supply
the substitution law and chain rule.  Degree-one cotangent classes then
construct structured square-zero extensions, and finite composites of their
underlying arrows produce the finite spectral infinitesimal substitutions
used later in the dissertation
\cite{LurieDAGIV,LurieHigherAlgebra}.

The cofree exponential differentiates inside the stable coefficient fibre,
whereas cotangent classes generate nonlinear infinitesimal geometry.  Their
common first-order interface is the cotangent classifier: on free
finite-projective models, the convolution derivative obtained from the
cofree calculus compares with the universal spectral K\"ahler derivation.
The same cotangent object then controls square-zero variation and reappears in
the reconstruction theory.

Finite infinitesimal substitutions provide the next structure needed by the
programme: they admit a Zariski semantics of their own.  The resulting
endpoint adjunctions are constructed first and only then recognized as a
differential-cohesive configuration.  Stabilizing that endpoint semantics
produces an exact-adjoint fracture pattern.  Its formal shape agrees with
familiar pullback constructions in differential cohomology, but the
comparison imports no independent curvature, period, cycle, or transfer
package
\cite{HopkinsSinger2005,BNV2016}.  At this stage the dissertation has
produced two differential mechanisms, their cotangent interface, and a
finite-infinitesimal geometry on which the later reconstruction and descent
problems can be posed.


\section{From infinitesimal data to reconstruction, descent, and jets}
\label{sec:intro-architecture}

Chapter~3 reverses the preceding direction and asks when first-order and
finite-infinitesimal data determine the nonlinear geometry from which they
arose.  It addresses this reconstruction problem through two approximation
coordinates, one measuring homotopical height and the other formal order at a
chosen centre.

Postnikov height resolves a connective spectral algebra into canonical
square-zero stages and recovers it as the inverse limit of its truncations.
Formal order is supplied instead by the canonical adic filtration of an
augmentation \(A\to B\).  Completion preserves the associated graded, and
the resulting linear layers are controlled by the shifted relative cotangent
complex:
\[
    \operatorname{Gr}
    \widehat{\operatorname{adic}}(A\to B)
    \simeq
    \operatorname{Sym}^{\mathrm{gr}}_B
    \bigl(L_{B/A}[-1]\langle1\rangle\bigr).
\]
The two coordinates answer different questions: Postnikov height
measures how much homotopical information remains, while the adic filtration
measures order of vanishing at the chosen centre.  This part of the
dissertation is developed in close contact with derived completion and formal
integration
\cite{BhattCompletionsDerivedDeRham,AntieauFiltrationsCohomologyIII2025,BrantnerMagidsonNuitenFormalIntegration}.

On a specified fixed-centre domain, the formal-order construction becomes a
reconstruction theorem.  Let \(B\) be a connective centre over a connective
base \(R\).  On connective almost finitely presented augmentations
\(A\to B\) that are surjective on \(\pi_0\) and fixed by canonical adic
completion, the shifted-cotangent/structured-square-zero adjunction is
comonadic.  Complete augmentations on this domain are reconstructed from
anchored shifted-cotangent data together with the coalgebra coherence generated
by
structured square-zero formation.  The bounded coherent sector
recovers the spectral Artinian test objects used in formal-moduli theory.

This fixed-centre result is not the only reconstruction mechanism in the
chapter.  A separate centre-free construction imposes Postnikov continuity
and nilpotent excision and produces a left-exact reconstruction localization.
Returning to a centre, split square-zero lifting yields an exact tangent
functor on eventually connective coefficients; for represented targets it is
represented by the cotangent complex, and the corresponding mapping-spectrum
functor extends canonically to all modules.  The crucial interface is not an
identification of the reconstruction comonad with the fibrewise differential
exponential.  Rather, the cotangent
classifier appears again as the common boundary object through which the
constructions can be compared.

Reconstruction asks when infinitesimal data determine nonlinear geometry.
Chapter~4 asks a different question: what remains when finite infinitesimal
variation is declared invisible?  Square-zero substitutions generate an
accessible left-exact de Rham localization, in the spirit of infinitesimal
and crystalline geometry
\cite{Ogus1975InfinitesimalSite,GaitsgoryRozenblyum2014Crystals,BeilinsonDrinfeld2004}.
For a relative context \(X\to S\), this produces the universal relative
de Rham reflection
\[
    X\longrightarrow X_{\mathrm{dR}/S}.
\]

Universal localization alone is not enough for the later descent theory,
because its unit need not be an effective epimorphism.  The unit is therefore
factored through its effective image,
\[
    X
    \xrightarrow{e_{X/S}}
    Q_{X/S}
    \xrightarrow{m_{X/S}}
    X_{\mathrm{dR}/S}.
\]
The two targets play different roles.  The object
\(X_{\mathrm{dR}/S}\) is the universal de Rham-local reflection, whereas
\(Q_{X/S}\) is the effective quotient actually presented by the generalized
points of \(X\).  The \v{C}ech nerve of
\(e_{X/S}\) is the effective de Rham relation, and this effective relation,
rather than the universal reflection by itself, is the object that supports
the subsequent descent and jet constructions.

After stabilizing the slice semantics, a distinguished additive scalar object
provides stable scalar functions on the simplices of this relation.  These
functions form a cosimplicial commutative algebra
\[
    C^\bullet_{\mathrm{dR}}(X/S),
\]
whose totalization computes stable scalar functions on the effective quotient.
The complete totalization filtration exposes intrinsic de Rham forms as its
normalized layers; the connecting morphisms supply a homotopy-coherent
differential, while the remaining filtration data encode multiplication,
higher closedness, exact-form lifts, and primitive data.  Restriction to the
first structured infinitesimal neighbourhood recovers the stable realization of
the cotangent classifier together with the universal derivation.  A separate
reduced-classical comparison probes a different boundary: it measures
convergence toward the reduced classical sector rather than first-order
variation.

Once the effective relation has been constructed, the remaining jet-theoretic
problem is to make that relation act coherently on objects varying over
\(X\).  Pullback along
\[
    e_{X/S}:X\longrightarrow Q_{X/S}
\]
has both dependent adjoints,
\[
    \Sigma_{e_{X/S}}
    \dashv
    e_{X/S}^*
    \dashv
    \Pi_{e_{X/S}}.
\]
The two composites encode the same effective infinitesimal relation in
monadic and comonadic form:
\[
    \mathsf{Disk}^{\mathrm{dR},\infty}_{X/S}
    =
    e_{X/S}^*\Sigma_{e_{X/S}},
    \qquad
    J^\infty_{X/S}
    =
    e_{X/S}^*\Pi_{e_{X/S}}.
\]
The first is the formal-disk monad and the second the infinite-jet comonad.
This is an intrinsically infinite-order construction; no global finite-order
jet tower is used in its definition.

Effective descent then identifies quotient descent, formal-disk action, and
jet coaction as three presentations of the same infinitesimal compatibility:
\[
    (\mathcal H_{\Sp})_{/Q_{X/S}}
    \simeq
    \operatorname{Alg}_{\mathsf{Disk}^{\mathrm{dR},\infty}_{X/S}}
    \bigl((\mathcal H_{\Sp})_{/X}\bigr)
    \simeq
    \operatorname{Coalg}_{J^\infty_{X/S}}
    \bigl((\mathcal H_{\Sp})_{/X}\bigr).
\]
Crystalline transport is the relation-theoretic form of this same effective
descent datum; it is not an independently chosen field on an object.  This
places the spectral construction in direct conversation with synthetic and
categorical jet formalisms
\cite{Kock1980,Marvan1986,KhavkineSchreiber2017,GiotopoulosKhavkineSatiSchreiber2026}.

The jet comonad now supplies the formal equation calculus.  Differential
operators are co-Kleisli morphisms out of \(J^\infty_{X/S}\), generalized
equations are morphisms into infinite-jet objects, and Cartan prolongation is
the universal repair that equips an ambient equation with compatible
crystalline transport.  A prolongation-complete generalized equation
canonically acquires the resulting jet-coalgebra and effective-descent
structure and hence determines an intrinsic formally integrable PDE.
Here formal integrability means infinitesimal descent: it does not assert
analytic existence, regularity, ellipticity, or functional-analytic
convergence of solutions.

The global jet construction is deliberately infinite-order.  On represented
affine diagonal augmentations, the adic filtration of Chapter~3 provides
genuine finite formal-order stages with cotangent-controlled associated
graded.  A global tower \(J^k_{X/S}\), however, would require a
base-change-coherent simplicial globalization of those affine stages.  No
such tower is assumed in the construction of \(J^\infty\) or in the
formal-integrability results.  The effective quotient \(Q_{X/S}\) is thus the
common geometric source of de Rham descent, formal-disk action, jet coaction, and crystalline transport,
which present the same underlying infinitesimal compatibility in different
forms.


\section{Main results and relation to existing work}
\label{sec:intro-thesis-results}

The thesis of the dissertation is that the indexed geometry of connective
spectral contexts and fully stable quasi-coherent coefficients contains
enough structure to support an infinitesimal and jet-theoretic foundation for
spectral variational geometry.  The principal results can be summarized as
follows.

\begin{enumerate}

    \item
    \textbf{Spectral contexts and stable coefficients form a
    base-change-compatible indexed doctrine.}
    Connective spectral schemes, stable quasi-coherent coefficient fibres,
    substitution, spectral Zariski descent, and relative-spectrum
    comprehension are organized in one indexed geometry.  This supplies the
    fixed semantic background in which nonlinear context extension and stable
    coefficient transport can be used together while remaining distinct.

    \item
    \textbf{Two distinct differential mechanisms are constructed and compared
    through cotangent data.}
    Each stable coefficient fibre carries a cofree
    cocommutative-coalgebra differential calculus, while the global cotangent
    classifier represents nonlinear first-order variation and constructs
    structured square-zero extensions.  The two mechanisms are not
    identified: they meet on the stated free finite-projective comparison
    domain through the universal spectral K\"ahler derivation.  Closing
    square-zero extensions under finite composition produces the
    infinitesimal test geometry, whose endpoint semantics yields differential
    cohesion and, after stabilization, an exact-adjoint fracture theorem and
    its subsequent BNV-type recognition.

    \item
    \textbf{Complete fixed-centre geometry is reconstructed comonadically from
    shifted-cotangent data on its theorem domain.}
    Postnikov stages relate homotopical approximation to finite infinitesimal
    substitutions, while the canonical adic filtration supplies formal order
    and its cotangent-controlled associated graded.  For a connective centre,
    connective almost finitely presented augmentations that are surjective on
    \(\pi_0\) and fixed by canonical adic completion are recovered
    comonadically from anchored shifted-cotangent coalgebra data.  A separate
    centre-free construction uses Postnikov continuity and nilpotent excision
    to define a left-exact reconstruction localization.  Tangent recovery
    returns these reconstruction mechanisms to the cotangent interface.

    \item
    \textbf{Finite-infinitesimal invariance produces a left-exact de Rham
    reflector and an effective descent relation.}
    Effective-image factorization separates the universal de Rham-local
    reflection from the quotient actually presented by generalized points.
    Stable scalar descent on the resulting relation produces intrinsic forms,
    a homotopy-coherent de Rham differential, and the associated multiplicative
    and higher descent data.  Formal \'etale descent and affine density then
    provide internal comprehension objects for the resulting form judgments.

    \item
    \textbf{Effective descent generates formal disks, infinite jets,
    crystalline transport, and a formal PDE calculus.}
    Dependent adjoints along the effective quotient produce the formal-disk
    monad and infinite-jet comonad.  Effective descent identifies quotient
    objects, disk actions, jet coactions, and crystalline transport as
    presentations of one infinitesimal compatibility structure.  The jet
    comonad then organizes differential operators, generalized equations,
    solution objects, Cartan prolongation, and the intrinsic
    formal-integrability criterion.

\end{enumerate}

These results use several established mathematical traditions, but at
different interfaces of the construction.  Spectral and derived algebraic
geometry provide the connective geometric setting, quasi-coherent
coefficients, and cotangent deformation theory
\cite{LurieDAGIV,LurieDAGIX,LurieHigherAlgebra,LurieSAG}.  Categorical logic
and linear/nonlinear semantics provide language for indexed substitution and
context extension
\cite{JacobsCategoricalLogic,Benton1995LNL}, while differential categories,
differential linear logic, and spectral cofree coalgebras provide comparison
targets for the fibrewise differential branch
\cite{BluteCockettSeely2006Differential,EhrhardRegnier2003Differential,Lemay2021Coderelictions,BachmannBurklund2026Coalgebras}.
Formal integration and derived deformation theory supply the nearest
comparison for the fixed-centre reconstruction problem
\cite{LurieDAGX,BrantnerMagidsonNuitenFormalIntegration}.  Infinitesimal
sites, crystals, and derived formal geometry provide the surrounding
language for de Rham localization and crystalline descent
\cite{Ogus1975InfinitesimalSite,GaitsgoryRozenblyum2014Crystals,GaitsgoryRozenblyumStudyII},
and synthetic and categorical approaches to jets and PDEs provide the
principal comparison for the final jet-theoretic calculus
\cite{Kock1980,Marvan1986,KhavkineSchreiber2017,GiotopoulosKhavkineSatiSchreiber2026}.

The novelty is not that each of these mathematical ingredients is new in
isolation.  It is that the relevant operations are constructed and
connected inside a single connective--stable indexed architecture, with the
passages between nonlinear geometry, stable coefficients, infinitesimal
variation, reconstruction, effective descent, and jets supplied by explicit
universal constructions rather than by adding independent compatibility
data.

The comparison with differential cohomology is correspondingly limited.
The stabilized endpoint fracture theorem has the same formal
exact-adjunction shape as sheaf-of-spectra models of differential cohomology
\cite{HopkinsSinger2005,BNV2016}, but the dissertation does not identify its
stable infinitesimal coefficients with a geometric differential-cohomology
theory.  No independent curvature, period, cycle, differential-form, or
transfer package is supplied by this comparison.  The phrase \emph{BNV-type}
is used only for this formal recognition, after the stabilized endpoint fracture theorem has already been established.


\section{Foundational conventions, scope, and organization}
\label{sec:intro-foundations-scope}

The main ambient convention is the distinction between connective geometry
and stable coefficients.  Coordinate algebras lie in
\[
    \CAlg(\Sp)^{\mathrm{cn}},
\]
while module and quasi-coherent coefficient categories retain their full
stable structure.  All categories are infinity-categories unless a truncation
is specified.  Mapping spaces are written
\(\Map_{\mathcal C}(X,Y)\), and mapping spectra in stable
infinity-categories are written
\(\map_{\mathcal C}(X,Y)\).

A second convention concerns higher coherence.  Structure supplied
functorially by a construction is retained through that construction:
adjunctions provide units, counits, and mates; Cartesian and coCartesian
transport provide substitution coherence; Kan extensions and descent provide
comparison and gluing maps; and effective descent provides the simplicial
coherence used in crystalline transport.  Later notation suppresses these
higher coherences once their source has been fixed.  By contrast, affine
charts, local trivializations, representatives of homotopy classes, finite
square-zero factorizations, local equations, and auxiliary centres introduced
inside proofs are proof witnesses and are not retained as additional
structure.

This distinction is particularly important for finite infinitesimal arrows.
Membership in the finite-composition closure of square-zero substitutions is
an intrinsic property, whereas the length of a chosen factorization is not an
intrinsic infinitesimal order.  Formal order is supplied separately by the
canonical adic filtration on the domain where that construction applies.
Conditions such as almost finite presentation, coherence, completeness,
Artinianity, formal smoothness, and formal \'etaleness likewise specify
theorem domains rather than additional fields of structure.

The differential equations constructed here are formal and
spectral-geometric.  Intrinsic formal integrability means effective descent
along the finite-infinitesimal de Rham relation, equivalently an object over
\(Q_{X/S}\).  It does not imply analytic existence or regularity of
solutions.  The explicitly physical part of the broader programme lies beyond
the present scope: Lagrangian field theory, Euler--Lagrange constructions,
variational bicomplexes, higher gauge fields, and their physical
interpretation are not claimed in this dissertation.

The five technical chapters follow the dependency order established above.
Chapters~1--2 construct and differentiate the indexed spectral doctrine.
Chapter~3 asks when the resulting infinitesimal data reconstruct nonlinear
geometry.  Chapters~4--5 instead pass to finite-infinitesimal invariance,
effective de Rham descent, and the resulting crystalline jet calculus.
Together they provide the infinitesimal and jet-theoretic foundation on which
the later variational programme is intended to build.

\section*{Note on use of AI}
Generative AI systems were used during the research and development of this
dissertation as tools for mathematical investigation, discussion, and
editorial refinement.  Responsibility for the mathematical content,
arguments, and conclusions rests with the author.


\endgroup
